\chapter{Two-Dimensional Continuum Elements}
\label{ch:2d_continuum}

\section{Plane Elasticity}

Two-dimensional continuum elements model thin plates or long prismatic
bodies~\cite{bathe2006finite,hughes2000fem}. Two assumptions are common:
\begin{description}
    \item[Plane stress] Thin plates loaded in their plane: $\sigma_z = \tau_{xz} = \tau_{yz} = 0$.
    \item[Plane strain] Long bodies with constrained out-of-plane strain: $\varepsilon_z = \gamma_{xz} = \gamma_{yz} = 0$.
\end{description}

\subsection{Stress and Strain Vectors}

For both cases, the in-plane state is described by:
\begin{equation}
    \vect{\sigma} = \begin{bmatrix}\sigma_x \\ \sigma_y \\ \tau_{xy}\end{bmatrix},
    \qquad
    \vect{\varepsilon} = \begin{bmatrix}\varepsilon_x \\ \varepsilon_y \\ \gamma_{xy}\end{bmatrix}
    \label{eq:2d_stress_strain}
\end{equation}

\subsection{Constitutive Matrix}

\textbf{Plane stress} ($\nu$ = Poisson's ratio, $E$ = Young's modulus):
\begin{equation}
    \matr{D}_{\text{ps}} = \frac{E}{1-\nu^2}
    \begin{bmatrix} 1 & \nu & 0 \\ \nu & 1 & 0 \\ 0 & 0 & \frac{1-\nu}{2} \end{bmatrix}
    \label{eq:D_planestress}
\end{equation}

\textbf{Plane strain:}
\begin{equation}
    \matr{D}_{\text{pe}} = \frac{E}{(1+\nu)(1-2\nu)}
    \begin{bmatrix} 1-\nu & \nu & 0 \\ \nu & 1-\nu & 0 \\ 0 & 0 & \frac{1-2\nu}{2} \end{bmatrix}
    \label{eq:D_planestrain}
\end{equation}

\section{Constant Strain Triangle (CST)}

The CST (or T3) is the simplest 2D element: 3 nodes, 2 DOFs per node (6 total),
with linear displacement fields and constant strains throughout.

\subsection{Shape Functions}

For a triangle with nodes $(x_1,y_1)$, $(x_2,y_2)$, $(x_3,y_3)$ and element
area $A_e$, the shape functions are the \emph{area coordinates}:
\begin{equation}
    N_i(x,y) = \frac{1}{2A_e}(a_i + b_i x + c_i y), \quad i = 1,2,3
    \label{eq:cst_shape}
\end{equation}

\begin{figure}[htbp]
\centering
\begin{tikzpicture}[>=Stealth, scale=1.2]
  %% Triangle vertices
  \coordinate (N1) at (0,0);
  \coordinate (N2) at (3.5,0);
  \coordinate (N3) at (1.1,2.2);
  %% Shaded area
  \fill[gray!20] (N1) -- (N2) -- (N3) -- cycle;
  \draw[thick] (N1) -- (N2) -- (N3) -- cycle;
  \node at (1.7,0.6) {$A_e$};
  %% Nodes — labels placed in the free quadrant (opposite the DOF arrows)
  \filldraw (N1) circle (3pt);
  \node[below left=2pt] at (N1) {1};
  \filldraw (N2) circle (3pt);
  \node[below right=2pt] at (N2) {2};
  \filldraw (N3) circle (3pt);
  \node[above left=2pt] at (N3) {3};
  %% DOF arrows: assumed-positive sense (along $+x$ and $+y$) at every node
  %% Node 1
  \draw[->,thick,blue!70!black] (N1) -- ++(0.85,0) node[below=1pt] {\small $u_1$};
  \draw[->,thick,blue!70!black] (N1) -- ++(0,0.85) node[left=1pt]  {\small $v_1$};
  %% Node 2
  \draw[->,thick,blue!70!black] (N2) -- ++(0.85,0) node[above=1pt] {\small $u_2$};
  \draw[->,thick,blue!70!black] (N2) -- ++(0,0.85) node[right=1pt] {\small $v_2$};
  %% Node 3
  \draw[->,thick,blue!70!black] (N3) -- ++(0.85,0) node[above=1pt] {\small $u_3$};
  \draw[->,thick,blue!70!black] (N3) -- ++(0,0.85) node[right=1pt] {\small $v_3$};
  %% Coordinate axes (well below triangle)
  \draw[->] (-1.5,-1.0) -- (-0.6,-1.0) node[right] {$x$};
  \draw[->] (-1.5,-1.0) -- (-1.5,-0.1) node[above] {$y$};
\end{tikzpicture}
\caption{CST element: 3 nodes, 2 DOFs per node ($u_i$, $v_i$), element area $A_e$.
DOF arrows show the assumed-positive directions along $+x$ and $+y$.}
\label{fig:cst_element}
\end{figure}

where:
\begin{align*}
    a_1 &= x_2 y_3 - x_3 y_2, \quad b_1 = y_2 - y_3, \quad c_1 = x_3 - x_2 \\
    2A_e &= \det\begin{bmatrix}1 & x_1 & y_1 \\ 1 & x_2 & y_2 \\ 1 & x_3 & y_3\end{bmatrix}
\end{align*}
(and cyclic permutations for $i=2,3$).

\subsection{Strain-Displacement Matrix}

\begin{equation}
    \vect{\varepsilon} = \matr{B}\, \vect{u}^e, \qquad
    \matr{B} = \frac{1}{2A_e}
    \begin{bmatrix}
        b_1 &  0  & b_2 &  0  & b_3 &  0  \\
         0  & c_1 &  0  & c_2 &  0  & c_3 \\
        c_1 & b_1 & c_2 & b_2 & c_3 & b_3
    \end{bmatrix}
    \label{eq:B_cst}
\end{equation}
Note that $\matr{B}$ is constant $\Rightarrow$ strain is constant over the CST element.

\subsection{Element Stiffness Matrix}

\begin{equation}
    \matr{k}^e = \int_{A_e} \matr{B}\transpose \matr{D} \matr{B} \,t \dd A
    = t\, A_e\, \matr{B}\transpose \matr{D} \matr{B}
    \label{eq:cst_stiffness}
\end{equation}
where $t$ is the element thickness.

\begin{importantbox}
The CST stiffness is evaluated in \emph{closed form} (no numerical integration
required) because $\matr{B}$ and $\matr{D}$ are both constant over the element:
\[
    \matr{k}^e = t\, A_e\, \matr{B}\transpose \matr{D} \matr{B} \in \mathbb{R}^{6\times6}
\]
\end{importantbox}


% =============================================================================
\subsection{Worked Example: Two-Element CST Mesh in Uniaxial Tension}
\label{sec:cst_example}
% =============================================================================

\begin{example}[Two-element CST plate under uniaxial tension]
\label{ex:cst_tension}

A $1\ \text{m} \times 1\ \text{m}$ square plate of thickness $t = 0.01\ \text{m}$
is modelled with two CST elements formed by splitting the square along its diagonal,
as shown in \cref{fig:cst_example}. Material properties: $E = 200\ \text{GPa}$,
$\nu = 0.3$, plane stress. A uniform tensile stress $\sigma_x = 100\ \text{MPa}$
is applied to the right edge, equivalent to a total force
$F = \sigma_x \cdot t \cdot 1\ \text{m} = 1\ \text{MN}$
split equally between nodes~2 and~3 as $F/2 = 500\ \text{kN}$ each.
The left-edge nodes are restrained: pin at node~1 ($u_1 = v_1 = 0$) and
horizontal roller at node~4 ($u_4 = 0$).

\begin{figure}[htbp]
\centering
\begin{tikzpicture}[>=Stealth, x=3.2cm, y=3.2cm,
    force/.style={->, ultra thick, red!75!black},
    disp/.style={->, thick, blue!70!black}]

  %% Square outline and diagonal
  \fill[gray!12] (0,0) -- (1,0) -- (1,1) -- (0,1) -- cycle;
  \draw[thick] (0,0) -- (1,0) -- (1,1) -- (0,1) -- cycle;
  \draw[thick] (0,0) -- (1,1);

  %% Element labels
  \node at (0.62,0.25) {\small\textcircled{\small 1}};
  \node at (0.25,0.70) {\small\textcircled{\small 2}};

  %% Nodes
  \filldraw (0,0) circle (2.5pt) node[below left=4pt] {\small\bfseries 1 $(0,0)$};
  \filldraw (1,0) circle (2.5pt) node[below right=2pt] {\small\bfseries 2 $(1,0)$};
  \filldraw (1,1) circle (2.5pt) node[above right=2pt] {\small\bfseries 3 $(1,1)$};
  \filldraw (0,1) circle (2.5pt) node[above right=2pt] {\small\bfseries 4 $(0,1)$};

  %% Pin support at node 1
  \draw[fill=gray!30,thick] (0,0) -- (-0.12,-0.18) -- (0.12,-0.18) -- cycle;
  \fill[pattern=north east lines,pattern color=gray!70]
       (-0.14,-0.18) rectangle (0.14,-0.25);
  \draw[thick] (-0.14,-0.18) -- (0.14,-0.18);

  %% Roller at node 4 (horizontal roller: free v, fixed u)
  \draw[fill=gray!30,thick] (0,1) -- (-0.12,0.82) -- (-0.12,1.18) -- cycle;
  \fill[pattern=north east lines,pattern color=gray!70]
       (-0.25,0.82) rectangle (-0.18,1.18);
  \draw[thick] (-0.18,0.82) -- (-0.18,1.18);

  %% Applied forces at nodes 2 and 3
  \draw[force] (1.0,0) -- (1.45,0)
       node[right=2pt] {$F/2 = 500\ \text{kN}$};
  \draw[force] (1.0,1) -- (1.45,1)
       node[right=2pt] {$F/2 = 500\ \text{kN}$};

  %% Dimensions
  \draw[<->,thin] (0,-0.35) -- (1,-0.35)
       node[fill=white,inner sep=1pt,midway] {$1\ \text{m}$};
  \draw[<->,thin] (-0.55,0) -- (-0.55,1)
       node[fill=white,inner sep=1pt,midway,rotate=90] {$1\ \text{m}$};

\end{tikzpicture}
\caption{Two-element CST mesh: 4 nodes, 2 triangular elements formed by the
diagonal from node~1 to node~3. Pin support at node~1, horizontal roller at
node~4, and $500\ \text{kN}$ horizontal force at nodes~2 and~3.}
\label{fig:cst_example}
\end{figure}

\medskip
\noindent\textbf{DOF ordering:}
$[u_1,v_1,u_2,v_2,u_3,v_3,u_4,v_4] = \text{DOFs}\ [1\ldots8]$.

\begin{table}[htbp]
\centering
\small
\caption{Node coordinates and element connectivity.}
\label{tab:cst_elements}
\begin{tabular}{@{}cccccc@{}}
\toprule
Elem & Nodes & Node coords (m) & $A_e$ (m$^2$) & Global DOFs \\
\midrule
1 & $1\to2\to3$ & $(0,0),(1,0),(1,1)$ & $0.5$ & 1,2,3,4,5,6 \\
2 & $1\to3\to4$ & $(0,0),(1,1),(0,1)$ & $0.5$ & 1,2,5,6,7,8 \\
\bottomrule
\end{tabular}
\end{table}


\medskip
\noindent\textbf{Step 1: Element geometry and B-matrices.}

For a CST with nodes at $(x_1,y_1)$, $(x_2,y_2)$, $(x_3,y_3)$, the
coefficients are $b_i = y_j - y_k$ and $c_i = x_k - x_j$
(cyclic, \cref{eq:cst_shape}), and $\matr{B}$ is constant (\cref{eq:B_cst}).

\smallskip
\noindent\textit{Element~1} (nodes $1\to2\to3$, coordinates $(0,0),(1,0),(1,1)$):
\[
    b_1 = -1,\quad b_2 = 1,\quad b_3 = 0,\qquad
    c_1 = 0,\quad c_2 = -1,\quad c_3 = 1,\qquad 2A_e = 1\ \text{m}^2
\]
\begin{equation}
    \matr{B}^{(1)} =
    \begin{bmatrix}
        -1 &  0 &  1 &  0 &  0 &  0 \\
         0 &  0 &  0 & -1 &  0 &  1 \\
         0 & -1 & -1 &  1 &  1 &  0
    \end{bmatrix}\ \text{m}^{-1}
    \label{eq:B1_cst_ex}
\end{equation}

\noindent\textit{Element~2} (nodes $1\to3\to4$, coordinates $(0,0),(1,1),(0,1)$):
\[
    b_1 = 0,\quad b_2 = 1,\quad b_3 = -1,\qquad
    c_1 = -1,\quad c_2 = 0,\quad c_3 = 1,\qquad 2A_e = 1\ \text{m}^2
\]
\begin{equation}
    \matr{B}^{(2)} =
    \begin{bmatrix}
         0 &  0 &  1 &  0 & -1 &  0 \\
         0 & -1 &  0 &  0 &  0 &  1 \\
        -1 &  0 &  0 &  1 &  1 & -1
    \end{bmatrix}\ \text{m}^{-1}
    \label{eq:B2_cst_ex}
\end{equation}

The two B-matrices differ because the diagonal splits the square asymmetrically:
element~1 has its right-angle corner at node~2, element~2 at node~4.

\medskip
\noindent\textbf{Step 2: Constitutive matrix.}

Plane stress with $E = 200\ \text{GPa}$, $\nu = 0.3$:
\begin{equation}
    \matr{D} = \frac{200\times10^9}{1-0.09}
    \begin{bmatrix}1 & 0.3 & 0 \\ 0.3 & 1 & 0 \\ 0 & 0 & 0.35\end{bmatrix}
    =
    \begin{bmatrix}
        2.198\times10^{11} & 6.593\times10^{10} & 0 \\
        6.593\times10^{10} & 2.198\times10^{11} & 0 \\
        0 & 0 & 7.692\times10^{10}
    \end{bmatrix}\ \text{Pa}
    \label{eq:D_cst_ex}
\end{equation}

\medskip
\noindent\textbf{Step 3: Element stiffness matrices.}

Using $\matr{k}^e = t\,A_e\,\matr{B}\transpose\matr{D}\matr{B}$ (\cref{eq:cst_stiffness})
with $t\,A_e = 0.01 \times 0.5 = 0.005\ \text{m}^3$:

\begin{equation}
    \matr{k}^{(1)} = 10^8 \times
    \begin{bmatrix}
        10.99 &  0     & -10.99 &  3.30  &  0     & -3.30  \\
         0    &  3.85  &   3.85 & -3.85  & -3.85  &  0     \\
       -10.99 &  3.85  &  14.84 & -7.14  & -3.85  &  3.30  \\
         3.30 & -3.85  &  -7.14 & 14.84  &  3.85  & -10.99 \\
         0    & -3.85  &  -3.85 &  3.85  &  3.85  &  0     \\
        -3.30 &  0     &   3.30 & -10.99 &  0     & 10.99
    \end{bmatrix}\ \tfrac{\text{N}}{\text{m}}
    \quad \text{(DOFs }u_1,v_1,u_2,v_2,u_3,v_3\text{)}
    \label{eq:ke1_cst}
\end{equation}

\begin{equation}
    \matr{k}^{(2)} = 10^8 \times
    \begin{bmatrix}
         3.85 &  0     &  0     & -3.85  & -3.85  &  3.85  \\
         0    & 10.99  & -3.30  &  0     &  3.30  & -10.99 \\
         0    & -3.30  & 10.99  &  0     & -10.99 &  3.30  \\
        -3.85 &  0     &  0     &  3.85  &  3.85  & -3.85  \\
        -3.85 &  3.30  & -10.99 &  3.85  & 14.84  & -7.14  \\
         3.85 & -10.99 &  3.30  & -3.85  & -7.14  & 14.84
    \end{bmatrix}\ \tfrac{\text{N}}{\text{m}}
    \quad \text{(DOFs }u_1,v_1,u_3,v_3,u_4,v_4\text{)}
    \label{eq:ke2_cst}
\end{equation}

\begin{remark}
$\matr{k}^{(1)} \neq \matr{k}^{(2)}$ even though both elements have the same area
and material. The two triangles are geometrically distinct (different orientations of
the hypotenuse) so their B-matrices differ, leading to different stiffness matrices.
\end{remark}

\medskip
\noindent\textbf{Step 4: Global assembly.}

The $8\times8$ global stiffness $\stiff$ is assembled from the two element
matrices. Element~1 contributes to DOFs $\{1,2,3,4,5,6\}$; element~2 to
DOFs $\{1,2,5,6,7,8\}$. DOFs $\{1,2,5,6\}$ (nodes~1 and~3) receive
contributions from both elements:

\begin{equation}
    \stiff = 10^8 \times
    \begin{bmatrix*}[r]
        14.84 &  0     & -10.99 &  3.30  &  0     & -7.14  & -3.85  &  3.85  \\
         0    & 14.84  &  3.85  & -3.85  & -7.14  &  0     &  3.30  & -10.99 \\
       -10.99 &  3.85  & 14.84  & -7.14  & -3.85  &  3.30  &  0     &  0     \\
         3.30 & -3.85  & -7.14  & 14.84  &  3.85  & -10.99 &  0     &  0     \\
         0    & -7.14  & -3.85  &  3.85  & 14.84  &  0     & -10.99 &  3.30  \\
        -7.14 &  0     &  3.30  & -10.99 &  0     & 14.84  &  3.85  & -3.85  \\
        -3.85 &  3.30  &  0     &  0     & -10.99 &  3.85  & 14.84  & -7.14  \\
         3.85 & -10.99 &  0     &  0     &  3.30  & -3.85  & -7.14  & 14.84
    \end{bmatrix*}\ \tfrac{\text{N}}{\text{m}}
    \label{eq:K_cst_global}
\end{equation}

Note that all diagonal entries equal $14.84 \times 10^8\ \text{N/m}$ --- a
consequence of the mesh symmetry.

\medskip
\noindent\textbf{Step 5: Load vector and boundary conditions.}

The global load vector has two non-zero entries:
$F_3 = F_5 = 500\,000\ \text{N}$ (forces at DOFs~3 and~5, i.e.\ $u_2$ and $u_3$).

Essential BCs: $u_1 = 0$ (DOF~1), $v_1 = 0$ (DOF~2), $u_4 = 0$ (DOF~7).
The five free DOFs are $\{3,4,5,6,8\}$ ($u_2, v_2, u_3, v_3, v_4$). The
reduced $5\times5$ system is:

\begin{equation}
    10^8 \times
    \begin{bmatrix}
        14.84 & -7.14  & -3.85  &  3.30  &  0     \\
        -7.14 & 14.84  &  3.85  & -10.99 &  0     \\
        -3.85 &  3.85  & 14.84  &  0     &  3.30  \\
         3.30 & -10.99 &  0     & 14.84  & -3.85  \\
         0    &  0     &  3.30  & -3.85  & 14.84
    \end{bmatrix}
    \begin{bmatrix} u_2 \\ v_2 \\ u_3 \\ v_3 \\ v_4 \end{bmatrix}
    =
    \begin{bmatrix} 500\,000 \\ 0 \\ 500\,000 \\ 0 \\ 0 \end{bmatrix}\ \text{N}
    \label{eq:Kff_cst}
\end{equation}

\medskip
\noindent\textbf{Step 6: Solution (MATLAB-verified, \Cref{lst:cst2d}).}

\begin{align*}
    u_2 &= +5.000\times10^{-4}\ \text{m}, &
    v_2 &= 0\ \text{m}, \\
    u_3 &= +5.000\times10^{-4}\ \text{m}, &
    v_3 &= -1.500\times10^{-4}\ \text{m}, \\
    u_4 &= 0\ \text{m}, &
    v_4 &= -1.500\times10^{-4}\ \text{m}.
\end{align*}

The right edge moves uniformly rightward ($u_2 = u_3 = 0.5\ \text{mm}$) with
no shear distortion. The top-right and top-left nodes contract vertically by
$0.15\ \text{mm}$ due to the Poisson effect.

\medskip
\noindent\textbf{Step 7: Stress recovery.}

For each element, the constant strain and stress vectors are:
\[
    \vect{\varepsilon}^e = \matr{B}^e\,\vect{u}^e, \qquad
    \vect{\sigma}^e = \matr{D}\,\vect{\varepsilon}^e .
\]

\noindent\textit{Element~1} ($\vect{u}^{(1)} = [0,0,5\!\times\!10^{-4},0,5\!\times\!10^{-4},-1.5\!\times\!10^{-4}]\transpose$ m):
\begin{align*}
    \vect{\varepsilon}^{(1)} &= \matr{B}^{(1)}\vect{u}^{(1)}
    = \begin{bmatrix}5\times10^{-4} \\ -1.5\times10^{-4} \\ 0\end{bmatrix}
    \begin{bmatrix}\text{m/m}\\\text{m/m}\\\text{rad}\end{bmatrix} \\[4pt]
    \vect{\sigma}^{(1)} &= \matr{D}\,\vect{\varepsilon}^{(1)}
    = \begin{bmatrix}100 \\ 0 \\ 0\end{bmatrix}\ \text{MPa}
\end{align*}

\noindent\textit{Element~2} ($\vect{u}^{(2)} = [0,0,5\!\times\!10^{-4},-1.5\!\times\!10^{-4},0,-1.5\!\times\!10^{-4}]\transpose$ m):
\[
    \vect{\varepsilon}^{(2)} = \vect{\varepsilon}^{(1)}, \qquad
    \vect{\sigma}^{(2)} = \vect{\sigma}^{(1)} =
    \begin{bmatrix}100 \\ 0 \\ 0\end{bmatrix}\ \text{MPa}
\]

Both elements give identical stresses: the CST mesh is in a state of uniform
$\sigma_x = 100\ \text{MPa}$, $\sigma_y = 0$, $\tau_{xy} = 0$.

\medskip
\noindent\textbf{Step 8: Analytical verification.}

The exact solution for a uniformly loaded elastic plate is:
\begin{align*}
    \sigma_x &= \frac{F}{t \cdot L} = \frac{1\times10^6}{0.01 \times 1}
              = 100\ \text{MPa}\ \checkmark \\
    \sigma_y &= 0\ \checkmark, \qquad \tau_{xy} = 0\ \checkmark \\[4pt]
    u(x,y) &= \frac{\sigma_x}{E}\,x = 5\times10^{-4}\,x\ \text{m}
              \;\Rightarrow\; u_2 = u_3 = 5\times10^{-4}\ \text{m}\ \checkmark \\
    v(x,y) &= -\frac{\nu\sigma_x}{E}\,y = -1.5\times10^{-4}\,y\ \text{m}
              \;\Rightarrow\; v_3 = v_4 = -1.5\times10^{-4}\ \text{m}\ \checkmark
\end{align*}

\begin{importantbox}
The CST recovers the \emph{exact} analytical solution for uniform stress
states. This is because the true displacement field ($u$ linear in $x$,
$v$ linear in $y$) lies within the CST polynomial space --- a polynomial
that the element can represent exactly. The FEM error is zero for any
problem whose exact solution is linear.
\end{importantbox}

The deformed mesh, magnified $\times200$, is shown in \cref{fig:cst_results}.

\begin{figure}[H]
\centering
\includegraphics[width=0.62\textwidth]{figures/cst_example_results}
\caption{CST mesh (solid) and deformed shape (dashed, magnified $\times200$).
The right edge moves uniformly rightward; the top nodes contract downward by
the Poisson effect. No shear distortion occurs.}
\label{fig:cst_results}
\end{figure}

\end{example}

\section{Isoparametric Quadrilateral Element (Q4)}

The 4-node quadrilateral (Q4) uses the same polynomials to interpolate both
geometry and displacement — this is the \emph{isoparametric} concept.

\subsection{Shape Functions in Natural Coordinates}

In the reference square $(\xi, \eta) \in [-1,1]^2$:
\begin{equation}
    N_i(\xi,\eta) = \frac{1}{4}(1 + \xi_i\xi)(1 + \eta_i\eta), \quad i=1,2,3,4
    \label{eq:q4_shape}
\end{equation}
where $(\xi_i, \eta_i)$ are the corner coordinates: $(-1,-1),(1,-1),(1,1),(-1,1)$.

\subsection{Jacobian Mapping}

The physical coordinates are mapped as:
\begin{equation}
    x = \sum_{i=1}^4 N_i\, x_i, \qquad y = \sum_{i=1}^4 N_i\, y_i
    \label{eq:q4_map}
\end{equation}
The Jacobian matrix $\matr{J}$ relates derivatives:
\begin{equation}
    \matr{J} = \begin{bmatrix} \pd{x}{\xi} & \pd{y}{\xi} \\ \pd{x}{\eta} & \pd{y}{\eta} \end{bmatrix}
    = \sum_{i=1}^4
    \begin{bmatrix} \pd{N_i}{\xi} \\ \pd{N_i}{\eta} \end{bmatrix}
    \begin{bmatrix} x_i & y_i \end{bmatrix}
    \label{eq:jacobian}
\end{equation}

Physical derivatives are recovered as:
\begin{equation}
    \begin{bmatrix} \pd{N_i}{x} \\ \pd{N_i}{y} \end{bmatrix}
    = \matr{J}\inv
    \begin{bmatrix} \pd{N_i}{\xi} \\ \pd{N_i}{\eta} \end{bmatrix}
\end{equation}

\subsection{Numerical Integration (Gauss Quadrature)}

The element stiffness integral is evaluated by $2\times2$ Gauss quadrature
(sufficient for Q4):
\begin{equation}
    \matr{k}^e = t \int_{-1}^{1}\int_{-1}^{1}
    \matr{B}(\xi,\eta)\transpose \matr{D}\, \matr{B}(\xi,\eta)\, |\matr{J}| \dd\xi\dd\eta
    \approx t \sum_{i=1}^{n_g} \sum_{j=1}^{n_g} w_i w_j\,
    \matr{B}(\xi_i,\eta_j)\transpose \matr{D}\, \matr{B}(\xi_i,\eta_j)\, |\matr{J}(\xi_i,\eta_j)|
    \label{eq:q4_stiffness}
\end{equation}
For $2\times2$ Gauss: $n_g=2$, $\xi_i = \pm 1/\sqrt{3}$, $w_i = 1$.

\begin{center}
\begin{tabular}{@{}ccc@{}}
    \toprule
    Order & Points per direction & Exact for polynomials up to \\
    \midrule
    1 & 1 & degree 1 \\
    2 & 2 & degree 3 \\
    3 & 3 & degree 5 \\
    \bottomrule
\end{tabular}
\end{center}

% =============================================================================
\subsection{Worked Example: Single Q4 Element (Cook Membrane)}
\label{sec:q4_example}
% =============================================================================

\begin{example}[Cook membrane --- single Q4 element under in-plane shear]
\label{ex:q4_cook}

The Cook membrane is a tapered quadrilateral plate used as a benchmark for
2D elements. A single Q4 element discretises the plate, as shown in
\cref{fig:q4_cook}. Material properties: $E = 70{,}000\ \text{N/mm}^2$,
$\nu = 0.3$, $t = 1\ \text{mm}$, plane stress. The left edge is clamped
(nodes~1 and~4). A total vertical force $P = 100\ \text{N}$ acts uniformly
on the right edge; the consistent nodal forces are $F/2 = 50\ \text{N}$
upward at each of nodes~2 and~3.

\begin{figure}[htbp]
\centering
\begin{tikzpicture}[>=Stealth, x=0.07cm, y=0.07cm,
    force/.style={->,ultra thick,red!75!black}]

  %% Plate outline
  \fill[gray!12]  (0,0) -- (48,44) -- (48,60) -- (0,44) -- cycle;
  \draw[thick]    (0,0) -- (48,44) -- (48,60) -- (0,44) -- cycle;

  %% Clamped left edge
  \fill[pattern=north east lines,pattern color=gray!70]
       (-5,0) rectangle (0,44);
  \draw[thick] (0,0) -- (0,44);

  %% Nodes
  \filldraw (0, 0) circle (3pt) node[left=4pt] {\small\bfseries 1 $(0,0)$};
  \filldraw (48,44) circle (3pt) node[right=4pt] {\small\bfseries 2 $(48,44)$};
  \filldraw (48,60) circle (3pt) node[right=4pt] {\small\bfseries 3 $(48,60)$};
  \filldraw (0, 44) circle (3pt) node[left=4pt] {\small\bfseries 4 $(0,44)$};

  %% Forces at nodes 2 and 3
  \draw[force] (48,44) -- ++(0,14) node[right=2pt,midway] {$50\ \text{N}$};
  \draw[force] (48,60) -- ++(0,14) node[right=2pt,midway] {$50\ \text{N}$};

  %% Dimensions
  \draw[<->,thin] (0,-8) -- (48,-8)
       node[fill=white,inner sep=1pt,midway] {$48\ \text{mm}$};
  \draw[<->,thin] (-10,0) -- (-10,44)
       node[fill=white,inner sep=1pt,midway,rotate=90] {$44\ \text{mm}$};
  \draw[<->,thin] (54,44) -- (54,60)
       node[fill=white,inner sep=1pt,midway,rotate=90] {$16\ \text{mm}$};

\end{tikzpicture}
\caption{Single Q4 element for the Cook membrane benchmark.
Node coordinates in mm. Left edge clamped; $50\ \text{N}$ vertical
nodal forces at nodes~2 and~3.}
\label{fig:q4_cook}
\end{figure}

\medskip
\noindent\textbf{DOF ordering:}
$[u_1,v_1,u_2,v_2,u_3,v_3,u_4,v_4] = \text{DOFs}\ [1\ldots8]$.\\
Fixed: $\{1,2,7,8\}$. Free: $\{3,4,5,6\}$ ($u_2,v_2,u_3,v_3$).

\medskip
\noindent\textbf{Step 1: Shape function natural derivatives.}

From \cref{eq:q4_shape}, the derivatives are:

\begin{table}[htbp]
\centering
\small
\caption{Natural derivatives of Q4 shape functions.}
\label{tab:q4_dN}
\begin{tabular}{@{}ccc@{}}
\toprule
Node $i$ & $\partial N_i/\partial\xi$ & $\partial N_i/\partial\eta$ \\
\midrule
1 & $-(1-\eta)/4$ & $-(1-\xi)/4$ \\
2 & $+(1-\eta)/4$ & $-(1+\xi)/4$ \\
3 & $+(1+\eta)/4$ & $+(1+\xi)/4$ \\
4 & $-(1+\eta)/4$ & $+(1-\xi)/4$ \\
\bottomrule
\end{tabular}
\end{table}

\medskip
\noindent\textbf{Step 2: Symbolic Jacobian.}

Substituting the node coordinates $(x_1,y_1)=(0,0)$,
$(x_2,y_2)=(48,44)$, $(x_3,y_3)=(48,60)$, $(x_4,y_4)=(0,44)$
into \cref{eq:jacobian}:
\begin{align}
    J_{11} &= \frac{\partial x}{\partial\xi}
    = \sum_i \frac{\partial N_i}{\partial\xi} x_i
    = \frac{-(1-\eta)}{4}\cdot0 + \frac{(1-\eta)}{4}\cdot48
      + \frac{(1+\eta)}{4}\cdot48 + \frac{-(1+\eta)}{4}\cdot0
    = 24 \notag\\
    J_{12} &= \frac{\partial y}{\partial\xi}
    = \frac{(1-\eta)}{4}\cdot44 + \frac{(1+\eta)}{4}\cdot60
      - \frac{(1+\eta)}{4}\cdot44
    = 15 - 7\eta \notag\\
    J_{21} &= \frac{\partial x}{\partial\eta}
    = \frac{-(1+\xi)}{4}\cdot48 + \frac{(1+\xi)}{4}\cdot48
    = 0 \notag\\
    J_{22} &= \frac{\partial y}{\partial\eta}
    = \frac{-(1+\xi)}{4}\cdot44 + \frac{(1+\xi)}{4}\cdot60
      + \frac{(1-\xi)}{4}\cdot44
    = 15 - 7\xi \notag
\end{align}

\begin{equation}
    \matr{J}(\xi,\eta) =
    \begin{bmatrix} 24 & 15-7\eta \\ 0 & 15-7\xi \end{bmatrix},
    \qquad
    \det|\matr{J}| = 24(15-7\xi)
    \label{eq:jacobian_cook}
\end{equation}

\begin{equation}
    \matr{J}^{-1} =
    \frac{1}{24(15-7\xi)}
    \begin{bmatrix} 15-7\xi & -(15-7\eta) \\ 0 & 24 \end{bmatrix}
    =
    \begin{bmatrix}
        \dfrac{1}{24} & \dfrac{-(15-7\eta)}{24(15-7\xi)} \\[8pt]
        0 & \dfrac{1}{15-7\xi}
    \end{bmatrix}
    \label{eq:Jinv_cook}
\end{equation}

\begin{remark}
$J_{21} = 0$ because the left and right edges are both vertical
($x_1 = x_4 = 0$, $x_2 = x_3 = 48\ \text{mm}$), so $x$ depends only
on $\xi$, not $\eta$. Consequently $\det|\matr{J}| = 24(15-7\xi)$
varies with $\xi$ only --- Gauss points at the same $\xi$ share the
same determinant regardless of $\eta$.
\end{remark}

\medskip
\noindent\textbf{Step 3: Gauss-point table.}

For $2\times2$ quadrature with $g = 1/\sqrt{3}$, $w_i = 1$:

\begin{table}[htbp]
\centering
\small
\caption{Jacobian data at the four Gauss points ($g=1/\sqrt{3}\approx0.5774$).
$J_{11}=24$ and $J_{21}=0$ at all points; $J^{-1}_{11}=1/24\approx0.04167$
and $J^{-1}_{21}=0$ at all points.}
\label{tab:q4_gauss}
\begin{tabular}{@{}cccccccc@{}}
\toprule
GP & $(\xi,\eta)$ & $J_{12}=15{-}7\eta$ & $J_{22}=15{-}7\xi$
   & $\det|\matr{J}|$ & $J^{-1}_{12}$ & $J^{-1}_{22}$ \\
\midrule
1 & $(-g,-g)$ & 19.04 & 19.04 & 457.0 & $-1/24$ & 0.05252 \\
2 & $(+g,-g)$ & 19.04 & 10.96 & 263.0 & $-0.07240$ & 0.09125 \\
3 & $(+g,+g)$ & 10.96 & 10.96 & 263.0 & $-1/24$ & 0.09125 \\
4 & $(-g,+g)$ & 10.96 & 19.04 & 457.0 & $-0.02398$ & 0.05252 \\
\bottomrule
\end{tabular}
\end{table}

\medskip
\noindent\textbf{Step 4: Strain-displacement matrix at GP1.}

The physical derivatives $\partial N_i/\partial x$ and $\partial N_i/\partial y$
follow from $[\partial N/\partial x,\; \partial N/\partial y]\transpose =
\matr{J}^{-1}[\partial N/\partial\xi,\; \partial N/\partial\eta]\transpose$.
At GP1 $(\xi\!=\!\eta\!=\!-g)$, substituting $J^{-1}_{12} = -1/24$
and $J^{-1}_{22} = 1/19.04$ (\cref{eq:Jinv_cook}):

\begin{align*}
    \frac{\partial N_i}{\partial x}
    &= \frac{1}{24}\frac{\partial N_i}{\partial\xi}
      -\frac{1}{24}\frac{\partial N_i}{\partial\eta},\\[4pt]
    \frac{\partial N_i}{\partial y}
    &= \frac{1}{19.04}\frac{\partial N_i}{\partial\eta}.
\end{align*}

Evaluating for all four nodes yields (in $10^{-3}\ \text{mm}^{-1}$):
\[
    \frac{\partial\vect{N}}{\partial x}
    = [0,\; 20.83,\; 0,\; -20.83],\qquad
    \frac{\partial\vect{N}}{\partial y}
    = [-20.71,\; -5.549,\; 5.549,\; 20.71]
\]

Assembling $\matr{B} = \matr{B}(\xi,\eta)$ at GP1:
\begin{equation}
    \matr{B}(\text{GP1}) = 10^{-3}\times
    \begin{bmatrix}
         0     & 0      & 20.83  & 0      & 0     & 0     & -20.83 & 0     \\
         0     & -20.71 & 0      & -5.549 & 0     & 5.549 & 0      & 20.71 \\
       -20.71  & 0      & -5.549 & 20.83  & 5.549 & 0     & 20.71  & -20.83
    \end{bmatrix}\ \text{mm}^{-1}
    \label{eq:B_q4_gp1}
\end{equation}

\begin{remark}
Unlike the CST where $\matr{B}$ is constant, the Q4 $\matr{B}(\xi,\eta)$
changes at every Gauss point because the physical derivatives depend on
$\matr{J}^{-1}(\xi,\eta)$, which varies over the element. Gauss quadrature
is therefore \emph{essential} --- no closed-form integral exists.
\end{remark}

\medskip
\noindent\textbf{Step 5: Constitutive matrix.}

\begin{equation}
    \matr{D} = \frac{70{,}000}{1-0.09}
    \begin{bmatrix}1&0.3&0\\0.3&1&0\\0&0&0.35\end{bmatrix}
    =
    \begin{bmatrix}
        76{,}923 & 23{,}077 & 0 \\
        23{,}077 & 76{,}923 & 0 \\
        0 & 0 & 26{,}923
    \end{bmatrix}\ \text{N/mm}^2
    \label{eq:D_q4_ex}
\end{equation}

\medskip
\noindent\textbf{Step 6: Element stiffness matrix.}

The $8\times8$ element stiffness is accumulated over the four Gauss points
(\cref{eq:q4_stiffness}), with all weights $w_i = 1$:
\[
    \matr{k}^e = t\sum_{p=1}^{4} \matr{B}_p\transpose\,\matr{D}\,\matr{B}_p\,
    \det|\matr{J}_p|
\]
Since $\matr{B}$ varies at each Gauss point, this sum must be computed
numerically (\Cref{lst:q4cook}). The MATLAB-verified result is:

\begin{equation}
    \matr{k}^e = 10^3\times
    \begin{bmatrix*}[r]
         14.33 &  -1.28 &   0.98 &  -9.95 &  -0.98 &  -3.51 & -14.33 &  14.74 \\
         -1.28 &  34.78 &  -8.02 &  19.75 &  -3.51 & -19.75 &  12.82 & -34.78 \\
          0.98 &  -8.02 &  85.76 & -34.66 & -37.68 &  21.20 & -49.06 &  21.49 \\
         -9.95 &  19.75 & -34.66 &  77.89 &  23.13 & -61.06 &  21.49 & -36.58 \\
         -0.98 &  -3.51 & -37.68 &  23.13 &  37.68 &  -9.66 &   0.98 &  -9.95 \\
         -3.51 & -19.75 &  21.20 & -61.06 &  -9.66 &  61.06 &  -8.02 &  19.75 \\
        -14.33 &  12.82 & -49.06 &  21.49 &   0.98 &  -8.02 &  62.40 & -26.28 \\
         14.74 & -34.78 &  21.49 & -36.58 &  -9.95 &  19.75 & -26.28 &  51.60
    \end{bmatrix*}\ \tfrac{\text{N}}{\text{mm}}
    \label{eq:ke_q4_cook}
\end{equation}

\medskip
\noindent\textbf{Step 7: Boundary conditions and reduced system.}

Clamping the left edge sets DOFs $\{1,2,7,8\} = 0$. The four free DOFs
$\{3,4,5,6\}$ ($u_2,v_2,u_3,v_3$) satisfy the $4\times4$ reduced system:

\begin{equation}
    10^3\times
    \begin{bmatrix}
         85.76 & -34.66 & -37.68 &  21.20 \\
        -34.66 &  77.89 &  23.13 & -61.06 \\
        -37.68 &  23.13 &  37.68 &  -9.66 \\
         21.20 & -61.06 &  -9.66 &  61.06
    \end{bmatrix}
    \begin{bmatrix} u_2 \\ v_2 \\ u_3 \\ v_3 \end{bmatrix}
    =
    \begin{bmatrix} 0 \\ 50 \\ 0 \\ 50 \end{bmatrix}\ \text{N}
    \label{eq:Kff_q4}
\end{equation}

\medskip
\noindent\textbf{Step 8: Solution (MATLAB-verified, \Cref{lst:q4cook}).}

\begin{align*}
    u_2 &= -9.3\times10^{-5}\ \text{mm}, &
    v_2 &= +8.210\times10^{-3}\ \text{mm}, \\
    u_3 &= -2.926\times10^{-3}\ \text{mm}, &
    v_3 &= +8.597\times10^{-3}\ \text{mm}.
\end{align*}

The dominant response is vertical (in the direction of loading): both right-edge
nodes deflect upward by approximately $8.4\ \mu\text{m}$. The small horizontal
displacements ($\sim\!3\ \mu\text{m}$) reflect the shear-induced distortion of
the tapered geometry. The deformed shape is shown in \cref{fig:q4_cook_results}.

\begin{figure}[H]
\centering
\includegraphics[width=0.62\textwidth]{figures/q4_cook_results}
\caption{Cook membrane: undeformed (solid) and deformed shape (dashed,
magnified $\times5000$). Both right-edge nodes deflect vertically;
the tapered geometry induces small lateral movement.}
\label{fig:q4_cook_results}
\end{figure}

\medskip
\noindent\textbf{Step 9: Stress recovery at Gauss points.}

The stresses $\vect{\sigma}^e = \matr{D}\,\matr{B}(\xi_p,\eta_p)\,\vect{u}^e$
are evaluated at each Gauss point:

\begin{table}[htbp]
\centering
\small
\caption{Stresses at the four Gauss points (N/mm$^2$).}
\label{tab:q4_stresses}
\begin{tabular}{@{}ccccc@{}}
\toprule
GP & $(\xi,\eta)$ & $\sigma_x$ & $\sigma_y$ & $\tau_{xy}$ \\
\midrule
1 & $(-g,-g)$ & $-0.100$ & $+0.121$ & $+4.181$ \\
2 & $(+g,-g)$ & $+5.436$ & $+2.608$ & $+1.608$ \\
3 & $(+g,+g)$ & $+0.173$ & $+1.029$ & $+1.860$ \\
4 & $(-g,+g)$ & $-3.128$ & $-0.788$ & $+4.327$ \\
\bottomrule
\end{tabular}
\end{table}

\begin{remark}
The stresses vary significantly across the element --- a key difference
from the CST, which yields constant stress. The large $\tau_{xy}$ values
at GPs~1 and~4 (left side of the element) reflect the shear dominance
near the clamped edge, where bending is most constrained. This stress
variation within the element is what the Q4 brings over the CST: it can
represent linearly varying strains, which better captures the true
behaviour of the Cook membrane under in-plane shear.
\end{remark}

\end{example}
